\documentclass[10pt,a4paper]{article}
\usepackage{amsmath,amssymb,amsthm,graphicx}
\usepackage{amsmath,amssymb,amsthm,hyperref,xfrac}
\providecommand{\ie}{i.e.\ }
\providecommand{\skpt}{\smallskip}
\usepackage[scr=rsfs,cal=euler]{mathalfa}
\multlinegap0pt

\hyphenation{begin decom-po-si-tion sequence}
\newenvironment{enumeratea}
{\bgroup\def\theenumi{\alph{enumi}}\begin{enumerate}}
{\end{enumerate}\egroup}

\newcommand{\ba}{{\boldsymbol a}}
\newcommand{\bb}{{\boldsymbol b}}
\newcommand{\bd}{{\boldsymbol d}}
\newcommand{\bu}{{\boldsymbol u}}
\newcommand{\bx}{{\boldsymbol x}}
\newcommand{\bmy}{{\boldsymbol y}}

\newcommand{\bD}{{\boldsymbol D}}
\newcommand{\bL}{{\boldsymbol L}}
\newcommand{\bP}{{\boldsymbol P}}
\newcommand{\bV}{{\boldsymbol V}}
\newcommand{\bX}{{\boldsymbol X}}
\newcommand{\bZ}{{\boldsymbol Z}}

\newcommand{\Psfrac}[2]{(\sfrac{#1}{#2})}
\newcommand{\psfrac}[2]{\sfrac{(#1)}{#2}}
\newcommand{\bigO}[1]{O(#1)}

\usepackage{mathtools}
\usepackage{bbm}

\newcommand{\NN}{\mathbb{N}}
\newcommand{\ZZ}{\mathbb{Z}}
\newcommand{\RR}{\mathbb{R}}
\newcommand{\CC}{\mathbb{C}}
\newcommand{\PP}{\mathbb{P}}
\newcommand{\EE}{\mathbb{E}}

\newcommand{\boldalpha}{\boldsymbol{\alpha}}
\newcommand{\bolddelta}{\boldsymbol{\delta}}
\newcommand{\boldrho}{\boldsymbol{\rho}}
\newcommand{\dee}{\,\mathrm{d}}
\newcommand{\one}{\mathbbm{1}}
\newcommand{\BC}{\mathcal{B}}
\newcommand{\EC}{\mathcal{E}}
\newcommand{\RS}{\mathscr{R}}
\newcommand{\prob}[1]{\PP[ #1 ]}
\newcommand{\bigoo}[1]{O\!\left(#1 \right)}
\newcommand{\abs}[1]{\left\lvert #1 \right\rvert}
\newcommand{\floor}[1]{\left\lfloor #1 \right\rfloor}
\newcommand{\riemint}[4]{\int_{#1}^{#2}\! #3 \, \mathrm{d}#4}
\newcommand{\maxnorm}[1]{\left\| #1 \right\|_{\infty}}

\renewcommand{\Re}{\textup{Re}}
\renewcommand{\Im}{\textup{Im}}

\newcommand{\dtv}{\mathrm{d}_{\mathrm{TV}}}
\newcommand{\dw}{\mathrm{d}_{\mathrm W}}

\let\epsilon\varepsilon

\newcommand{\pextra}{P_{\mathtt{extra}}}
\newcommand{\parratia}{P_{\mathtt{Arratia}}}

\newcommand{\simplex}{\Delta^{k-1}}

\newcommand{\bg}{\big}
\newcommand{\bgg}{\Big}
\newcommand{\bggg}{\bigg}

\newcommand{\Rnt}{R_{\mathrm{NT}}}
\newcommand{\Rprob}{R_{\mathrm{PD}}}
\newcommand{\ECnt}{\EC_{\mathrm{NT}}}
\newcommand{\ECprob}{\EC_{\mathrm{PD}}}

\newcommand{\Bnt}{\mathcal{B}_{\mathrm{NT}}}
\newcommand{\Bprob}{\mathcal{B}_{\mathrm{PD}}}

\DeclareMathOperator{\vard}{Var}
\DeclareMathOperator{\Dir}{Dir}

\theoremstyle{plain}
\newtheorem{cor}{Corollary}[section]
\newtheorem{lem}[cor]{Lemma}
\newtheorem{prop}[cor]{Proposition}
\newtheorem{conj}[cor]{Conjecture}
\newtheorem{theo}{Theorem}

\theoremstyle{definition}
\newtheorem{example}{Example}
\newtheorem{deff}[cor]{Definition}
\newtheorem{rem}{Remark}[section]
\newtheorem*{rems}{Remark}
\newtheorem*{remss}{Remarks}

\makeatletter
\ams@newcommand{\multiint}[1]{\DOTSI\protect\MultiIntegral{#1}}
\renewcommand{\MultiIntegral}[1]{%
\edef\ints@c{\noexpand\intop
\ifnum#1=\z@\noexpand\intdots@\else\noexpand\intkern@\fi
\replicate{#1-2}{\noexpand\intop\noexpand\intkern@}%
\noexpand\intop
\noexpand\ilimits@
}%
\futurelet\@let@token\ints@a
}
\makeatother
\ExplSyntaxOn
\cs_new:Npn \replicate #1 #2 { \prg_replicate:nn { #1 } { #2 } }
\ExplSyntaxOff
\usepackage{cleveref}
\crefname{originalTheorem}{original Theorem}{original Theorems}
\Crefname{originalTheorem}{Original Theorem}{Original Theorems}
\makeatletter
\newlabel{thm:fact into k parts}{{2}{source}{}{sourcepointer-1}{}}
\newlabel{thm:fact into k parts@cref}{{[originalTheorem][0][]2}{[1][0][]source}{}{}{}}
\makeatother
\title{Uniform Dirichlet boundary bound with a dependent prime-rounding error radius}
\author{Tony Haddad, Dimitris Koukoulopoulos}
\date{doi:10.5802/jep.317; CC BY4.0}
\begin{document}
\maketitle
\noindent\textit{Editorial note.} The visible joint-model setup names Z\_i exactly as the colored mass vector at original1123-1125, with independent iid colors independent of V. It states x>=2 and positive simplex parameters. Theta\_x uses that same V; there is no independence assumption between Z\_i and Theta\_x. This is a setup clarification only; the retained author proof is unchanged.\par\medskip
\noindent\textit{Editorial note.} The event has fixed center u and random radius2Theta\_x/log x; Theta\_x depends on the same Poisson-Dirichlet sequence that defines Z\_i. No independence between Z\_i and Theta\_x is assumed. The integer coupling N\_x,M\_x and its old public expected-distance theorem are not needed for this probability-only result.\par\medskip
\noindent\textit{Editorial note.} At1473 the beta integral is printed with reciprocal Gamma product. The preceding Euler identity1454-1455 and the explicit identity1611 uniquely give sin(pi alpha)*Gamma(alpha)*Gamma(1-alpha)=pi. Preserve the original erroneous intermediate display and disclose its reciprocal slip.\par\medskip
\noindent\textit{Editorial note.} In the dyadic display1524, the preceding construction1521 and the goal1529 specify Theta\_x-prime>2\textasciicircum{}m; the displayed unprimed Theta\_x and u\_i in place of fixed u are notation slips. The rest of the proof consistently controls Theta\_x-prime.\par\medskip
\noindent\textit{Editorial note.} The summand at1541 uses V\_i while the sum and both surrounding lines1540/1542 run over j. The original sum1510 and dyadic count definition1535 uniquely specify V\_j. The author text is retained, not silently altered.\par\medskip
\noindent\textit{Editorial note.} The Campbell formula is used only for nonnegative summands, where its displayed integrability condition applies. This task does not assert the printed unrestricted signed-function version in1817-1823.\par\medskip
\noindent\textit{Editorial note.} The h codomain at469 is written strictly positive although the step sum is0 below lambda\_0; the formula472 and explicit observation482 determine the nonnegative rounding function used in r(t)=|h(t)-t|. No numerical value of h is changed.\par\medskip
\noindent\textit{Editorial scope.} One dependent-radius probability bound. The standalone beta integral estimate, gamma/Dirichlet coloring identity and dyadic second-moment lemma are not counted separately. The full integer-factorization Dirichlet theorem is not claimed prepared by this card.\par\medskip
\noindent\textit{Allowed standard prerequisites (editorial).}\begin{itemize}
\item Poisson-Dirichlet gamma-Poisson representation. Descending points of a Poisson process with intensity e\textasciicircum{}(-t)dt/t, normalized by their sum, have parameter-one Poisson-Dirichlet law.. Represent the joint colored components and the radius on the same point process.
\item Poisson coloring, Campbell/Laplace formula and multivariate Mecke equation. Split into independent color processes and remove/reinsert two atoms, preserving their contribution to the random normalization.
\item Gamma normalization, Euler reflection and elementary beta integrals. The gamma sum computation and uniform beta-window estimate are printed locally; only standard identities and gamma-to-Dirichlet normalization are free.
\item Classical prime number theorem and strong Mertens estimate. Justify the explicit rounding envelope r(t)<=c min(t,t\textasciicircum{}(-2)) at451-478, not the correlated event estimate.
\end{itemize}\medskip
\ifdefined\hypertarget\hypertarget{sourcepointer-1}{}\fi
\noindent\textit{Source reference.} Original Theorem 2: sources/arratia-dirichlet-jep-2025-317/source0.tex, source lines 349--359. The section introduces the probabilistic model of this later theorem. The boundary estimate is proved independently from its colored-PD setup; the integer-factorization theorem is neither an input nor separately counted here.\par\medskip
\paragraph{Joint probability model (editorial setup from the author construction).}
Fix $x\ge 2$, an integer $k\ge 2$, and positive parameters
$\alpha_1,\ldots,\alpha_k$ with $\sum_{i=1}^k\alpha_i=1$.
Let $\mathbf V=(V_j)_{j\ge1}$ have the parameter-one Poisson--Dirichlet law.
Let $(C_j)_{j\ge1}$ be independent and identically distributed colors,
independent of $\mathbf V$, with $\mathbb P(C_j=i)=\alpha_i$.
Define, on this same probability space,
\[
 Z_i=\sum_{j:\,C_j=i}V_j\qquad (1\le i\le k).
\]
Use the author's $\Theta_x$ formed from this very sequence $\mathbf V$.
Thus $Z_i$ and $\Theta_x$ are generally dependent.
This setup names the colored masses already defined in the author's
Poisson--Dirichlet coloring proposition; it adds no integer-factorization coupling.

There are many ways to define the Poisson--Dirichlet distribution. One of the most intuitive ones involves a ``stick-breaking'' process that we will use throughout the paper. We~start by sampling a sequence of i.i.d.~random variables $(U_i)_{i \ge 1}$ that are all uniformly distributed in $[0, 1]$. We~then define the sequence $(L_i)_{i \ge 1}$ in the following way:
\[
L_1 \coloneqq U_1 \quad \text{and}\quad L_j \coloneqq U_j \prod_{i=1}^{j-1} (1-U_i) \quad \text{for $j \ge 2$.}
\]
The distribution of the process $\bL = (L_1, L_2, \ldots)$ is called the \textit{GEM distribution} (of~parameter 1), named after Griffiths, Engen and McCloskey (see \cite[Ch.\,41]{JKB} for a discussion on the history of this distribution). Lastly, we~sort the components of $\bL$ in non-increasing order to create $\bV = (V_1, V_2, \ldots)$. The distribution of this process is the \textit{Poisson--Dirichlet distribution} (of parameter 1).\footnote{The GEM and Poisson--Dirichlet distributions have more general definitions involving typically a parameter $\theta$. In~the rest of the paper, we~will not be mentioning the parameter since we will always work with $\theta = 1$.} We note that both $\sum_{i \ge 1} L_i$ and $\sum_{i \ge 1} V_i$ are equal to $1$ almost surely.

Let $\simplex$ be the set of $k$-tuples $\boldalpha=(\alpha_1, \ldots, \alpha_k) \in \RR^k_{\ge 0}$ satisfying $\alpha_1+\cdots+\alpha_k=1$. We~also need to define a special class of functions:

Starting with the pioneering work of Deshouillers, Dress, and Tenenbaum \cite{DDT79}, many different people have studied the distribution of $\bD_{f, x}$ for various choices of $f \in \mathcal F_k(\boldalpha)$, and they proved that it converges on a logarithmic scale to a certain {\it Dirichlet distribution}. Recall that if $\boldalpha \in \simplex$ satisfies $\alpha_i > 0$ for all $i$, we~say that a $\simplex$-valued random vector $\bZ$ follows the Dirichlet distribution $\Dir(\boldalpha)$ if
\[
\PP\bg[Z_i \le u_i\ \, \forall i<k\bg] = F_{\boldalpha}(\bu) \coloneqq \prod_{i=1}^{k} \Gamma(\alpha_i)^{-1} \mathop{\int\cdots\int}\limits_{\substack{0 \le t_i \le u_i \ \forall i < k \\ t_1 + \cdots + t_{k-1} \le 1}} \prod_{i=1}^k t_i^{\alpha_i - 1} \dee t_1\cdots \dee t_{k-1}.
\]

Let $(\lambda_j)_{j \ge 0}$ be the increasing sequence of positive real numbers defined by $\lambda_0 \coloneqq e^{-\gamma}$ and
\[
\lambda_j \coloneqq \exp\!\bgg(\!-\gamma+ \sum_{i \le j} \frac{1}{v_iq_i}\bgg)\quad \text{for $j \ge 1$},
\]
with $\gamma$ being the Euler-Mascheroni constant and $q_j = p_j^{v_j}$ being the $j$\textsuperscript{th} smallest prime power, \ie $(q_j)_{j \ge 1}$ is the sequence $2, 3, 2^2, 5, 7, 2^3, 3^2,\dots$ Note that
\begin{equation}
\label{eq:lambda_j}
\lambda_j = \log q_j +O(1/(\log q_j)^2)
\end{equation}
by a stronger version of Mertens's estimate (Proposition \ref{prop:mertens}).

Moreover, we~have
\begin{equation}
\label{eq:q_j}
q_{j+1} = q_j+O(q_j/(\log q_j)^3).
\end{equation}
Indeed, the prime number theorem (Proposition \ref{pnt}) implies that for some sufficiently large constant $C$, we~have $\theta(q_j+Cq_j/(\log q_j)^3) -\theta(q_j)>0$, whence $q_{j+1}\le q_j+Cq_j/(\log q_j)^3$, as required.

Next, we~define the step-function $h:\RR_{>0} \to \RR_{>0}$ by
\begin{equation}
\label{eq:h(t) dfn}
h(t) \coloneqq \sum_{j \ge 1} (\log q_j) \cdot \one_{\lambda_{j-1} < t \le \lambda_j}.
\end{equation}
In particular, \eqref{eq:lambda_j} and \eqref{eq:q_j} imply that, if $r(t) \coloneqq |h(t) - t|$, then
\begin{equation}
\label{eq:h(t) estimate}
r(t) \ll \min\{t, t^{-2}\} \quad \text{for all $t > 0$.}
\end{equation}

\begin{equation}
\label{eq:def of Theta_x}
\Theta_x \coloneqq \sum_{i \ge 1} r(V_i\log x).
\end{equation}

\[
\ECprob(\simplex_{\mathbf{u}})^c\subseteq \bigcup_{\substack{i< k \\ u_i\neq 1}} \Bprob^{(i)}(u_i),
\quad\text{where}\quad
\Bprob^{(i)}(u) \coloneqq \bgg\{ \bg| Z_i - u \bg|\le \sfrac{2\Theta_x}{\log x} \bgg\},
\]

\begin{prop}[The prime number theorem]
\label{pnt}
For $x \ge 2$, we~have
\[
\theta(x) = x + \bigO{\sfrac{x}{(\log x)^3}}
\]
and
\[
\psi(x) = x+\bigO{\sfrac{x}{(\log x)^3}}.
\]
\end{prop}

We shall also need a stronger version of Mertens's estimate, which in its classical form, yields an estimate for the product $\prod_{p \le x} (1-1/p)$.

\begin{prop}[Strong Mertens's estimate]
\label{prop:mertens}
For $x \ge 2$, we~have
\[
\sum_{p^k \le x} \frac{1}{kp^k} = \log_2 x + \gamma + \bigO{\sfrac{1}{(\log x)^3}}.
\]
\end{prop}

\begin{proof}
Using the prime number theorem (Proposition \ref{pnt}), there exists a constant $c_0$ satisfying
\[
\sum_{p^k \le x} \frac{1}{kp^k} = \int_{2^-}^x \frac{\dee \psi(t)}{t\log t} = \log_2 x + c_0 + \bigO{\sfrac{1}{(\log x)^3}}.
\]
for $x \ge 2$. For the proof of $\gamma = c_0$, see the proof of \cite[Th.\,2.7]{MV07}.
\end{proof}

\subsection{Poisson Processes}

The following definition and propositions are borrowed from Kingman's book on Poisson processes \cite{Kingman93}. Let $(S, \mathcal S)$ be a measurable space with~$S$ being a subset of $\RR^d$ for some $d \ge 1$. A \textit{Poisson process} on a state space $S$ with mean measure $\mu$ is a random countable subset $\Pi \subseteq S$ such that:
\begin{itemize}
\item for any disjoint measurable subsets $A_1, \ldots, A_n$ of $S$, the random variables \hbox{$\#(\Pi \cap A_1)$, \dots, $\#(\Pi \cap A_n)$} are independent;

\item the random variable $\#(\Pi \cap A)$ is a Poisson random variable of parameter $\mu(A)$ for any $A\subset S$ measurable.
\end{itemize}
We can see $\Pi$ as an element of the measurable space $(\Omega_S, \mathcal F_S)$ where $\Omega_S$ is the set of countable subsets of $S$ and $\mathcal F_S$ is the smallest $\sigma$-algebra for which the map $\Pi \mapsto \#\{\Pi \cap\nobreak B\}$ is measurable for all $B \in \mathcal S$. If~$\mu$ has no atoms, meaning no singleton with positive probability, and is $\sigma$-finite, meaning that $S$ is a countable union of measurable sets with finite measure, then a Poisson process with mean measure $\mu$ always exists (see \cite{Kingman93}, the existence theorem in section 2.5 for a proof). If~$\mu$ is absolutely continuous with respect to the Lebesgue measure, then the function \hbox{$\lambda \colon S \to \RR_{\ge 0}$} such that $\mu(A) = \riemint{A}{}{\lambda(x)}{x}$ for all measurable subsets $A \subseteq S$ is called the \textit{intensity} of the Poisson process. Here are a few important propositions about Poisson processes that we will use in the paper and proved in \cite{Kingman93}:

\begin{prop}[Mapping theorem, {\cite[\S2.3]{Kingman93}}]
\label{prel:map}
If $\Pi$ is a Poisson process with mean measure $\mu$ on $S$, and $f:S \to T$ is a measurable function such that $\mu^*(B) \coloneqq \mu(f^{-1}(B))$ has no atoms, then $f(\Pi)$ is a Poisson process on $T$ with measure $\mu^*$.
\end{prop}

\begin{prop}[Coloring theorem, {\cite[\S5.1]{Kingman93}}]
\label{prel:color}
If $\Pi$ is a Poisson process with mean measure $\mu$ on $S$, and the points are randomly colored with $k$ colors such that the probability of a point receiving the color $i$ is $p_i$, and such that the color of a point is independent of different points and of the position of the point. Let $\Pi_i$ be the subset of~$\Pi$ with color $i$. Then all the $\Pi_i$ are independent Poisson processes with mean measures $\mu_i = p_i\mu$.
\end{prop}

\begin{prop}[Campbell's theorem, {\cite[\S3.2]{Kingman93}}]
\label{prel:campbell}
Let $\Pi$ be any Poisson process on $S$ with mean measure $\mu$. Let $f\colon S \to \RR$ be a measurable function. Then
\[
\Sigma = \sum_{X \in \Pi} f(X)
\]
is absolutely convergent almost surely if, and only if,
\[
\riemint{S}{}{\min\{f, 1\}}{\mu} < \infty.
\]
If this condition holds, then
\[
\EE[e^{s\Sigma}] = \exp\left(\riemint{S}{}{(e^{sf} - 1)}{\mu}\right)
\]
for any complex $s$ for which the integral converges. Moreover,
\[
\EE[\Sigma] = \riemint{S}{}{f}{\mu}
\]
if the integral converges. In~the case where it converges, we~also have
\[
\vard[\Sigma] = \riemint{S}{}{f^2}{\mu}.
\]
\end{prop}

Many probabilities or expectations that involve Poisson processes in this paper can be reformulated as
\[
\EE \sum_{X \in \Pi} f(\Pi\setminus \{X\}, X).
\]
In these cases, there is a generalization of the formula for $\EE[\Sigma]$ in Campbell's theorem, called the Mecke equation, allowing us to compute these objects:

\begin{prop}[Mecke equation, {\cite[Th.\,4.5]{LastPenrose18}}]
\label{prel:mecke}
Let $\Pi$ be a Poisson process on $S$ with a $\sigma$-finite mean measure $\mu$, and let $f: \Omega_S \times S^k \to [0, \infty)$ be measurable. Then we have
\begin{multline*}
\EE \sum_{\substack{X_1, \ldots, X_k \\ \text{all distinct}}} f(\Pi\setminus \{X_1, \ldots, X_k\}; X_1, \ldots, X_k) \\[-10pt]
= \int_S \cdots \int_S \EE[f(\Pi; x_1, \ldots, x_k)] \, \mathrm d\mu(x_1)\, \cdots\, \mathrm d\mu(x_k).
\end{multline*}
\end{prop}

\begin{lem}
\label{lem:bound of prob of ECprob}
For $i\in\{1,2,\dots,k-1\}$ and $u\in(0,1)$, we~have
\[
\PP\bg[\Bprob^{(i)}(u) \bg] \ll \frac{1}{(1+ u\log x)^{1-\alpha_i}(1+(1-u)\log x)^{\alpha_i}}.
\]
\end{lem}

\section{Theorem \ref{thm:fact into k parts} for the probabilistic model}\label{sec:Donnelly-Tavare}

With the following probabilistic version of Theorem \ref{thm:fact into k parts} given by Donnelly and Tavaré in 1987 \cite[\S 3]{DonnellyTavare87}, we~see how a Poisson--Dirichlet process relates to the Dirichlet distribution. We~repeat their proof since they use a Poisson process that will come in handy later.

\begin{prop}[Donnelly--Tavaré \cite{DonnellyTavare87}]
\label{prop:probabilistic dirichlet law}
Let $\bV = (V_1, V_2, \ldots)$ be a Poisson--Dirichlet process, let $\boldalpha = (\alpha_1, \ldots, \alpha_k) \in \simplex$ such that $\alpha_i > 0$ for all $1 \le i \le k$, and let $(C_i)_{i \ge 1}$ be a sequence of i.i.d.~random variables (also independent of $\bV$) with $\PP[C_i = j] = \alpha_j$ for all $i \ge 1$. Then the random vector
\[
\Bigl(\sum_{i:\ C_i = 1} V_i, \ldots, \sum_{i:\ C_i = k} V_i\Bigr)
\]
is distributed according to $\Dir(\boldalpha)$.
\end{prop}

\begin{proof}
Let $X_1 > X_2 > \cdots$ be the points of a Poisson process on $\RR_{>0}$ with intensity $\Psfrac{e^{-x}}{x} \, \mathrm dx$, and let $S \coloneqq \sum_{i \ge 1} X_i$. For each $i \ge 1$, Let $V_i \coloneqq \sfrac{X_i}{S}$. We~then know that $(V_1, V_2, \ldots)$ follows the Poisson--Dirichlet distribution \cite[Th.\,2.2]{Feng10}. With the coloring theorem (Proposition \ref{prel:color}), the point processes
\[
\Pi_m = \{X_i : C_i = m\}
\]
form independent Poisson processes of intensity $\alpha_m\cdot \Psfrac{e^{-x}}{x} \, \mathrm dx$ for all $m = 1, \ldots, k$. We~then let
\[
S_m \coloneqq \sum_{X \in \Pi_m} X
\]
for $m=1,2,\dots,k$. With Campbell's theorem (Proposition \ref{prel:campbell}), we~compute the moment generating function
\[
\log \EE[e^{sS_m}] = \alpha_m\cdot \riemint{0}{\infty}{\frac{e^{sx}-1}{x}\cdot e^{-x}}{x} = -\alpha_m\log(1-s)
\]
for $\Re(s) < 1$. This coincides with the moment generating function of $\text{Gamma}(\alpha_m, 1)$ distribution. We~thus deduce that the vector
\[
\Bigl(\sum_{i:\ C_i = 1} V_i, \ldots, \sum_{i:\ C_i = k} V_i \Bigr) = \Bigl(\frac{S_1}{\sum_{m=1}^k S_m}, \ldots, \frac{S_k}{\sum_{m=1}^k S_m} \Bigr)
\]
follows the distribution $\text{Dir}(\boldalpha)$ (see \cite[Ch.\,49, pp.\,485-487]{KotzBal00} for a proof).
\end{proof}

\section{Proof of Lemma \ref{lem:bound of prob of ECprob} using probabilistic techniques}\label{sec:boundary event Prob}

Since $\bZ$ follows the distribution $\text{Dir}(\boldalpha)$, each component $Z_i$ follows a \hbox{$\text{Beta}(\alpha_i,\mkern-3mu 1 \!-\! \alpha_i)$} distribution, meaning that if $[a,b]\subseteq[0,1]$, then
\begin{equation}
\label{eq:Zi distribution}
\begin{aligned}
\PP\bg[Z_i\in[a,b] \bg] &= \frac{1}{\Gamma(\alpha_i)\Gamma(1-\alpha_i)}\int_a^b \frac{\dee t}{t^{1-\alpha_i} (1-t)^{\alpha_i}}\\
&= \frac{\sin(\pi\alpha_i)}{\pi} \int_a^b \frac{\dee t}{t^{1-\alpha_i} (1-t)^{\alpha_i}},
\end{aligned}
\end{equation}
by Euler's reflection formula. For this reason, we~need the following preliminary estimate.

\begin{lem}\label{lem:beta estimate}
Uniformly for $u\in[0,1]$, $\alpha \in (0, 1)$ and $\delta > 0$, we~have
\[
\sin(\pi\alpha) \int_{[u-\delta,u+\delta]\cap[0,1]} \frac{\dee t}{t^{1-\alpha}(1-t)^\alpha} \ll \frac{\delta}{(u+\delta)^{1-\alpha}(1-u+\delta)^\alpha}.
\]
In particular, if $\delta\ge1/\log x$, then
\[
\sin(\pi\alpha) \int_{[u-\delta,u+\delta]\cap[0,1]} \frac{\dee t}{t^{1-\alpha}(1-t)^\alpha} \ll \frac{\delta\log x}{(1+u\log x)^{1-\alpha}(1+(1-u)\log x)^\alpha}.
\]
\end{lem}

\begin{proof} Note that both sides of the claimed inequality are invariant under the change of variables $(\alpha,u)\mapsto (1-\alpha,1-u)$. Hence, we~may assume without loss of generality that $u\in[0,1/2]$. In~addition, observe that
\[
\sin(\pi\alpha)\int_0^1 \frac{\dee t}{t^{1-\alpha}(1-t)^\alpha} = \frac{\sin(\pi\alpha)}{\Gamma(\alpha)\Gamma(1-\alpha)} = \pi.
\]
Hence the lemma is trivially true if $\delta > 1/4$.

We have thus reduced the proof to the case when $\alpha\in(0,1)$, $u\in[0,1/2]$ and $\delta\in[0,1/4]$. In~particular, $u+\delta\le 3/4$, so that $1-t\in[1/4,1]$ for all $t\in[0,u+\delta]$. It~thus suffices to prove that
\begin{equation}
\label{eq:beta bound - new goal}
\sin(\pi\alpha) \int_{[u-\delta,u+\delta]\cap[0,1]} \frac{\dee t}{t^{1-\alpha}} \ll \frac{\delta}{(u+\delta)^{1-\alpha}}.
\end{equation}

Assume first that $\delta\le u/2$. We~then have $t\ge u/2$ whenever $t\in[u-\delta,u+\delta]$. Using also the trivial bound $\sin(\pi\alpha)\le1$, we~conclude that
\[
\sin(\pi\alpha) \int_{[u-\delta,u+\delta]\cap[0,1]} \frac{\dee t}{t^{1-\alpha}} \le \int_{[u-\delta,u+\delta]\cap[0,1]} \frac{\dee t}{(u/2)^{1-\alpha}} \ll \frac{\delta}{(u+\delta)^{1-\alpha}}.
\]
This proves the lemma in this case.

Finally, assume that $\delta\ge u/2$. We~then use that
\begin{align*}
\sin(\pi\alpha) \int_{[u-\delta,u+\delta]\cap[0,1]} \frac{\dee t}{t^{1-\alpha}}
&\le \sin(\pi\alpha)\int_0^{3\delta} \frac{\dee t}{t^{1-\alpha}} =\frac{\sin(\pi \alpha)}{\alpha} \cdot (3\delta)^\alpha \\
&\ll \delta^\alpha \asymp \frac{\delta}{(u+\delta)^{1-\alpha}}.
\end{align*}
This completes the proof of the lemma in all cases.
\end{proof}

Let us now show Lemma \ref{lem:bound of prob of ECprob}. For simplicity of notation, let us fix $i\in\{1,2,\dots,k\}$ and $u\in(0,1)$, and let us set $\alpha \coloneqq \alpha_i$ and
\[
\Delta \coloneqq \frac{1}{(1+ u\log x)^{1-\alpha}(1+(1-u)\log x)^{\alpha}}.
\]
Thus, our goal is to show that
\begin{equation}
\label{eq:Bprob rewrite}
\PP\bgg[ \bg| Z_i - u \bg|\le \sfrac{2\Theta_x}{\log x}\bgg] \ll \Delta.
\end{equation}

Recall that
\[
\Theta_x=\sum_{j\ge1}r(V_j\log x),
\]
and that there exists an absolute constant $c\ge1$ such that $r(y)\le c \min\{y,y^{-2}\}$ for all $y>0$ (see \eqref{eq:h(t) estimate}). Since there are at most $10$ indices $j$ such that $V_j\ge 0.1$, we~find
\[
\Theta_x\le \Theta_x'+10c,\quad\text{where}\quad \Theta_x'\coloneqq \sum_{j\ge1:\ V_j<0.1}r(V_j\log x).
\]

Now, using Lemma \ref{lem:beta estimate} and relation \eqref{eq:Zi distribution}, we~have
\[
\PP\bgg[ \bg| Z_i - u \bg|\le \sfrac{200c}{\log x}\bgg] \ll \Delta.
\]
On the other hand, if $\frac{200c}{\log x}<|Z_i-u|\le \frac{2\Theta_x}{\log x}\le \frac{20c+2\Theta_x'}{\log x}$, then we must have $\Theta_x'>90c$. In~particular, there exists $m\in\ZZ$ such that $2^m>40c$ and $\Theta_x'\in(2^m,2^{m+1}]$, whence $|Z_i-u|\le \frac{20c + 2^{m+2}}{\log x}$. We~thus conclude that
\[
\PP\bgg[ \bg| Z_i - u \bg|\le \frac{20c+2\Theta_x'}{\log x}\bgg] \le \sum_{m\in\ZZ:\ 2^m>40c}
\PP\bgg[ \bg| Z_i - u_i \bg|\le \frac{20c+2^{m+2}}{\log x},\ \Theta_x>2^m\bgg] + O(\Delta).
\]
Therefore Lemma \ref{lem:bound of prob of ECprob} will follow if we can prove that
\begin{equation}
\label{eq:lem Bprob reduction 1}
\PP\bgg[ \bg| Z_i - u \bg|\le \frac{20c+4\kappa}{\log x},\ \Theta_x'>\kappa \bgg] \ll \frac{\Delta}{\kappa}
\qquad\text{for all}\ \kappa\ge 40c.
\end{equation}

We shall make a further reduction. If~we set
\[
G(\lambda)\coloneqq \#\bg\{j\ge1: V_j\log x\in (\lambda,2\lambda] \bg\},
\]
then we have
\[
\begin{split}
\Theta_x' =\sum_{j\ge1:\ V_j<0.1} r(V_j\log x)
&\le c\sum_{j\ge1:\ V_j<0.1}\min\bg\{V_i\log x,(V_i\log x)^{-2}\bg\} \\
&\le c\sum_{\substack{m\ge0 \\ 2^m\le 0.1\log x}} \frac{G(2^m)}{4^m} + c \sum_{m<0} 2^{m+1} G(2^m) \\
&\le 5c \max_{\substack{m\ge0 \\ 2^m\le 0.1\log x}} \bgg(\frac{G(2^m)}{2^{3m/2}} \bgg) + 5c \max_{m<0}\bg(2^{m/2} G(2^m)\bg),
\end{split}
\]
since $\sum_{m\ge0}2^{-m/2} <5$ and $\sum_{m<0}2^{1+m/2}< 5$. We~thus find
\begin{multline*}
\PP\bgg[ \bg| Z_i - u \bg|\le \frac{20c+4\kappa}{\log x},\ \Theta_x'>\kappa \bgg] \\
\le \sum_{\substack{m\ge0 \\ 2^m\le 0.1\log x}} \!\!\! \PP\bgg[ \bg| Z_i - u \bg|\le \frac{20c+4\kappa}{\log x},\ G(2^m)> \frac{2^{3m/2} \kappa}{10c} \bgg] \\
+\sum_{m<0} \PP\bgg[ \bg| Z_i - u \bg|\le \frac{20c+4\kappa}{\log x},\ G(2^m)> \frac{\kappa 2^{-m/2}}{10c} \bgg].
\end{multline*}
Note that in both sums, we~have $G(2^m)\ge 4$, since $\kappa\ge 40c$. Hence, Markov's inequality implies
\begin{multline*}
\PP\bgg[ \bg| Z_i - u \bg|\le \frac{20c+4\kappa}{\log x},\ \Theta_x'>\kappa \bgg] \\
\le \sum_{\substack{m\ge0 \\ 2^m\le 0.1\log x}} \frac{100c^2}{\kappa^22^{3m}}
\EE\bgg[ \one_{| Z_i - u |\le \frac{20c+4\kappa}{\log x} } \cdot \one_{G(2^m) \ge4} \cdot G(2^m)^2\bgg] \\
+\sum_{m<0} \frac{100c^22^m}{\kappa^2}
\EE\bgg[ \one_{| Z_i - u |\le \frac{20c+4\kappa}{\log x} }\cdot \one_{G(2^m) \ge4}\cdot G(2^m)^2\bgg].
\end{multline*}

This reduces \eqref{eq:lem Bprob reduction 1}, and thus Lemma \ref{lem:bound of prob of ECprob}, to proving the following estimate:

\begin{lem}
\label{lem:Bprob reduction 2} Uniformly for $\mu\ge1$ and $\lambda\in(0,0.1\log x]$, we~have
\[
\EE\bgg[ \one_{| Z_i - u |\le \sfrac{\mu}{\log x} } \cdot \one_{G(\lambda) \ge4} \cdot G(\lambda)^2 \bgg]
\ll (\lambda+\mu)\cdot \Delta.
\]
\end{lem}

\begin{proof} Let us call $E$ the quantity we seek to bound from above.
Consider two independent Poisson processes $(A_j)_{j \ge 1}$ with intensity $\alpha\Psfrac{e^{-x}}{x} \dee x$ and $(B_j)_{j \ge 1}$ with intensity $(1-\alpha)\Psfrac{e^{-x}}{x} \dee x$. Hence, the union of the two processes is a new Poisson process of intensity $\Psfrac{e^{-x}}{x}\dee x$.
The sum $S_A \coloneqq \sum_{j\ge1} A_j$ has distribution $\text{Gamma}(\alpha, 1)$ and the sum $S_B \coloneqq \sum_{j\ge1} B_j$ has distribution $\text{Gamma}(1-\alpha, 1)$. If~we also set $S\coloneqq S_A+S_B$, then, as we saw in the proof of Proposition \ref{prop:probabilistic dirichlet law}, we~have
\[
E
=\EE\bgg[ \one_{| S_A/S- u |\le \sfrac{\mu}{\log x} } \cdot \one_{G_A(\lambda)+G_B(\lambda)\ge 4}\cdot \bg( G_A(\lambda)+G_B(\lambda) \bg)^2 \bgg],
\]
where $G_A(\lambda) \!=\! \sum_{j\ge1} \one_{A_j\log x\in(\lambda S,2\lambda S]}$ and $G_B(\lambda)\!=\! \sum_{j\ge1} \one_{B_j\log x\in(\lambda S,2\lambda S]}$. We~know $G_A(\lambda)+G_B(\lambda)\ge4$. Hence, if $G_A(\lambda)\ge G_B(\lambda)$, then we must have \hbox{$G_A(\lambda)\ge 2$} and $G_A(\lambda)\le G_A(\lambda)^2/2$, whence $G_A(\lambda) \le 2\mathop{\sum\sum}_{k>j\ge1} \one_{A_k\log x,\, A_j\log x\in (\lambda S,2\lambda S]}$. So, we~find
\[
\begin{split}
\bg( G_A(\lambda)+G_B(\lambda) \bg)^2 \le 4G_A(\lambda)^2
&= 4 G_A(\lambda) + 8 \mathop{\sum\sum}_{k>j\ge1} \one_{A_k\log x,\, A_j\log x\in (\lambda S,2\lambda S]} \\
&\le 16 \mathop{\sum\sum}_{k>j\ge1} \one_{A_k\log x,\, A_j\log x\in (\lambda S,2\lambda S]}.
\end{split}
\]
The analogous inequality also holds when $G_B(\lambda)\ge G_A(\lambda)$, with the roles of $A$ and~$B$ reversed. We~conclude that
\[
E
\le 16 (E_A+E_B),
\]
where
\[
E_A \coloneqq \EE\bgg[\textstyle \one_{|S_A/S - u |\le \sfrac{\mu}{\log x} } \mathop{\sum\sum}_{k>j\ge1} \one_{A_k\log x,\, A_j\log x\in (\lambda S,2\lambda S]} \bgg]
\]
and $E_B$ is defined analogously. Using Mecke's equation (cf. Proposition \ref{prel:mecke}), we~find
\[
\begin{split}
E_A &= \alpha^2\iint\limits_{2a_1>a_2>a_1>0}
\EE\bgg[\textstyle \one_{| \frac{a_1+a_2+S_A}{a_1+a_2+S} - u |\le \frac{\mu}{\log x} } \prod_{j=1}^2 \one_{a_j\log x \in (\lambda (a_1+a_2+S),2\lambda (a_1+a_2+S)]} \bgg] \\[-10pt]
&\hspace*{9.5cm} \times \frac{e^{-a_1-a_2}}{a_1a_2}\dee a_1\dee a_2 \\[6pt]
&= \frac{\alpha^2}{\Gamma(\alpha)\Gamma(1-\alpha)} \mkern-50mu \iiiint\limits_{
\substack{
2a_1>a_2>a_1>0,\ s_1,s_2>0 \\
| \frac{a_1+a_2+s_1}{a_1+a_2+s_1+s_2} - u |\le \frac{\mu}{\log x} \\
a_j\log x \in (\lambda (a_1+a_2+s_1+s_2),2\lambda (a_1+a_2+s_1+s_2)] \ (j=1,2)
} }
\mkern-50mu\frac{e^{-a_1-a_2-s_1-s_2}}
{a_1a_2s_1^{1-\alpha}s_2^\alpha}\dee a_1\dee a_2 \dee s_1\dee s_2,
\end{split}
\]
where we used that the $a_j$'s lie in the same dyadic interval to deduce that $a_2<2a_1$. We~make the change of variables $t=s_1/(s_1+s_2)$ and $s=s_1+s_2$. Since $\lambda/\log x\le 0.1$ and $a_1<a_2<2a_1$, the conditions $a_j\log x \in (\lambda (a_1+a_2+s),2\lambda (a_1+a_2+s)]$ for $j=1,2$ imply that $a_j\log x \in(\lambda s,5\lambda s]$. Knowing also that $ |\frac{a_1+a_2+s_1}{a_1+a_2+s} - u |\le \frac{\mu}{\log x}$, we~find $|t-u|=|\frac{s_1}{s_1+s_2} - u|\le \frac{10\lambda+\mu}{\log x}$. Finally, we~use Euler's reflection formula to write $\Gamma(\alpha)\Gamma(1-\alpha)=\pi/\sin(\pi\alpha)$. We~conclude that
\[
E_A \le \frac{\sin(\pi \alpha)}{\pi}
\iiiint\limits_{
\substack{
a_2>a_1>0,\ s>0,\ t\in[0,1] \\
| t-u|\le (10\lambda+\mu)/\log x\\
a_j\log x \in (\lambda s, 5\lambda s] \ (j=1,2)
} }
\frac{e^{-a_1-a_2-s}}
{a_1a_2t^{1-\alpha}(1-t)^\alpha}\dee a_1\dee a_2 \dee s\dee t.
\]
For every fixed value of $s$, the integral over $a_1$ and $a_2$ is $\le(\log 5)^2$, since $\int_w^{5w}\frac{e^{-a}}{a}\dee a \le \log5$ for any $w>0$. We~also have $\int_0^\infty e^{-s}\dee s=1$. We~thus conclude that
\[
E_A\le \frac{(\log 5)^2\sin(\pi \alpha) }{\pi} \int\limits_{\substack{t\in[0,1] \\ |t-u|\le (10\lambda+\mu)/\log x}} \frac{\dee t}{t^{1-\alpha}(1-t)^\alpha}.
\]
Using Lemma \ref{lem:beta estimate} shows that $E_A\ll (\lambda+\mu)\Delta$. The same estimate holds for $E_B$ too, thus completing the proof of the lemma.
\end{proof}

\begin{thebibliography}{99}
\bibitem{Arratia98} R. Arratia - “On the central role of scale invariant Poisson processes on (0,∞)”, in Microsurveys in discrete probability (Princeton, NJ, 1997), DIMACS Ser. Discrete Math. Theoret. Comput. Sci., vol. 41, American Mathematical Society, Providence, RI, 1998, p. 21-41
\bibitem{Arratia02} R. Arratia - “On the amount of dependence in the prime factorization of a uniform random integer”, in Contemporary combinatorics, Bolyai Soc. Math. Stud., vol. 10, János Bolyai Math. Soc., Budapest, 2002, p. 29-91
\bibitem{ABT06} R. Arratia, A. D. Barbour \& S. Tavaré - “A tale of three couplings: Poisson-Dirichlet and GEM approximations for random permutations”, Combin. Probab. Comput. 15 (2006) no. 1-2, p. 31-62
\bibitem{BareikisManstavicius07} G. Bareikis \& E. Manstavičius - “On the DDT theorem”, Acta Arith. 126 (2007) no. 2, p. 155-168
\bibitem{Bill72} P. Billingsley - “On the distribution of large prime divisors”, Period. Math. Hungar. 2 (1972), p. 283-289
\bibitem{delaBretecheTenenbaum16} R. de la Bretèche \& G. Tenenbaum - “Sur les processus arithmétiques liés aux diviseurs”, Adv. in Appl. Probab. 48 (2016) no. A, p. 63-76
\bibitem{DDT79} J.-M. Deshouillers, F. Dress \& G. Tenenbaum - “Lois de répartition des diviseurs. I”, Acta Arith. 34 (1979) no. 4, p. 273-285 (loose errata)
\bibitem{DonnellyTavare87} P. Donnelly \& S. Tavaré - “The population genealogy of the infinitely-many neutral alleles model”, J. Math. Biol. 25 (1987) no. 4, p. 381-391
\bibitem{Feng10} S. Feng - The Poisson-Dirichlet distribution and related topics. Models and asymptotic behaviors, Probability and its Appl., Springer, Heidelberg, 2010
\bibitem{JKB} N. L. Johnson, S. Kotz \& N. Balakrishnan - Discrete multivariate distributions, Wiley Series in Probab. and Stat.: Applied Probab. and Stat., John Wiley \& Sons, Inc., New York, 1997
\bibitem{Kingman93} J. F. C. Kingman - Poisson processes, Oxford Studies in Probab., vol. 3, The Clarendon Press, Oxford University Press, New York, 1993
\bibitem{KotzBal00} S. Kotz, N. Balakrishnan \& N. L. Johnson - Continuous multivariate distributions. Vol. 1. Models and applications, Wiley Series in Probab. and Stat.: Applied Probab. and Stat., Wiley-Interscience, New York, 2000
\bibitem{Koukoulo19} D. Koukoulopoulos - The distribution of prime numbers, Graduate Studies in Math., vol. 203, American Mathematical Society, Providence, RI, 2019
\bibitem{LastPenrose18} G. Last \& M. Penrose - Lectures on the Poisson process, Institute of Math. Stat. Textbooks, vol. 7, Cambridge University Press, Cambridge, 2018
\bibitem{Ken22} S.-K. Leung - “Dirichlet law for factorisation of integers, polynomials and permutations”, Math. Proc. Cambridge Philos. Soc. 175 (2023) no. 3, p. 649-676
\bibitem{KenWeb24} S.-K. Leung - “Dirichlet law with arbitrary parameters for factorizations of integers” (2025), Retrieved on 24 August 2025 from https://sites.google.com/view/sunkaileung/miscellaneous
\bibitem{MV07} H. L. Montgomery \& R. C. Vaughan - Multiplicative number theory. I. Classical theory, Cambridge Studies in Advanced Math., vol. 97, Cambridge University Press, Cambridge, 2007
\bibitem{NyandwiSmati13} S. Nyandwi \& A. Smati - “Distribution laws of pairs of divisors”, Integers 13 (2013), article ID A13, 13 pages
\bibitem{Tenenbaum00} G. Tenenbaum - “A rate estimate in Billingsley’s theorem for the size distribution of large prime factors”, Q. J. Math. 51 (2000) no. 3, p. 385-403
\end{thebibliography}
\end{document}
